\section{引言}

旅行商问题（Traveling Salesman Problem, TSP）是组合优化领域最为经典且最具挑战性的 NP-hard 问题之一。给定 $n$ 个城市及它们两两之间的距离，TSP 要求找到一条经过每座城市恰好一次并最终返回起点的最短闭合路径。尽管问题陈述极为简洁，其解空间的规模却随城市数量呈阶乘级增长（$(n-1)!/2$ 条可能路径），使得在 $n$ 较大时穷举搜索完全不现实。在工程实践中，TSP 具有广泛的应用背景——从 PCB 钻孔的最短路径规划、物流配送的路径优化到集成电路的 VLSI 布线，无不涉及 TSP 或其变体的求解需求。

进化算法（Evolutionary Algorithm, EA）作为一类模拟自然选择与遗传变异机制的随机搜索方法，不依赖问题的梯度信息或特定解析结构，已被证明能够在大规模 TSP 实例上有效逼近最优解。然而，单一种群的进化算法在处理较高维度的 TSP 实例时面临一个根本性的困境——早熟收敛（premature convergence）。在进化后期，种群中的所有个体趋于同质化，遗传多样性的枯竭使得变异和交叉算子难以再产生具有新颖结构的子代，导致算法陷入局部最优而无法进一步突破。这一现象的根源在于，有限规模的种群无法在整个搜索空间中同时维持充分的多样性和收敛精度——随着进化推进，遗传漂变（genetic drift）不可逆地消融了种群中曾经存在的遗传变异。

并行进化算法（Parallel Evolutionary Algorithm, PEA）通过引入多个协同进化的子种群，为破解这一困境提供了天然且有效的路径。在并行架构下，每个子种群在各自的计算单元上独立演化，在不同的搜索区域进行探索；同时，通过定期的迁移（migration），子种群之间有选择地交换个体，使得每个孤岛既能从他者的优秀探索成果中受益，又不至于因过度频繁的基因交流而丧失自身的探索独立性。Li 与 Hu 在发表于《Soft Computing》2014 年的论文《Global Migration Strategy with Moving Colony for Hierarchical Distributed Evolutionary Algorithms》中，基于这一思想提出了系统的分层分布式进化算法框架。该框架的核心创新包括两项：第一，将迁移分为局部（组内）和全局（组间）两个层次，以截然不同的频率和粒度运行，构建起真正的层级结构；第二，提出了"移动种群"（Moving Colony）的全局迁移策略，将全局迁移的对象从单个个体提升至整个子种群，使得全局迁移可以完全不依赖物理通信即可完成，且在单次迁移中为接收组注入的信息量远超于传统的个体迁移。

本报告以该论文为理论基础，使用 MPI（Message Passing Interface）并行编程模型，将原有的串行 TSP 进化算法完整并行化，实现了论文中描述的全部三种算法——基础的环形分布式进化算法（Ring DEA）、环形个体迁移的分层分布式进化算法（Ring Individual HDEA）以及基于移动种群的环形分层分布式进化算法（Ring Colony HDEA，亦称 HdEA-MP）。三种算法均在严格遵循论文参数设置的条件（64 子种群、8×8 组结构、Inver-over 算子、映射算子、临界速度机制、变异率衰减、20:1 局部-全局迁移比、100 全局轮次终止条件等）下进行 10 次独立运行，并基于 t 检验对三种算法的求解质量和稳定性进行了系统的定量比较与因果分析。
